% ============================================================
% 3. Method: Formal Model, Architecture, and Theorems
% ============================================================

\section{Method}
\label{sec:method}

We first define the formal model (\S\ref{sec:formal}), then present
the BIC architecture (\S\ref{sec:architecture}) and state our four
main theorems (\S\ref{sec:theorems}).  Full proofs are in
Appendix~\ref{sec:proofs}.

% ============================================================
\subsection{Formal Model}
\label{sec:formal}

\begin{definition}[Agent Swarm]
\label{def:swarm}
An \emph{agent swarm} is a rooted tree $\Tree = (V, E, \sigma, \delta)$
where:
\begin{itemize}
  \item $V$ is a (potentially unbounded) set of \emph{agent nodes},
        with root $v_0 \in V$.
  \item $E \subseteq V \times V$ are directed parent$\to$child edges.
  \item $\sigma: V \times \States \to V^*$ is the \emph{spawn function};
        executing agent $v$ with state $s$ may spawn children
        $\sigma(v,s) = (v_1, \dots, v_k)$.
  \item $\delta: V \times \States \to \States$ is the \emph{transition
        function}; executing agent $v$ in state $s$ produces a new
        state $\delta(v, s)$.
\end{itemize}
The \emph{state space} $\States$ is the set of all possible agent
states (e.g., JSON-serializable dictionaries).
\end{definition}

\begin{definition}[Bounded Infinity Cache (BIC)]
\label{def:bic}
A BIC is a tuple $\BIC = (\slots, h, \summarize, \info)$ where:
\begin{itemize}
  \item $\slots \in \Nat^+$ is the \emph{cache capacity}---the maximum
        number of agent states stored simultaneously.
  \item $h: V \to [\slots]$ is the \emph{slot mapping}: given an agent
        $v$ at tree position $(p_1, \dots, p_d)$ (sequence of child
        indices from root), define
        \begin{equation}
          h(v) = \cantor(p_1, \dots, p_d) \bmod \slots
          \label{eq:slot}
        \end{equation}
        where $\cantor: \Nat^* \to \Nat$ is the generalized Cantor
        pairing function (iterated left-fold of the standard Cantor
        pair $\pi(a,b) = \frac{(a+b)(a+b+1)}{2} + b$).
  \item $\summarize: \States \times \States^k \to \States$ is the
        \emph{summarization function}: given a parent state and $k$
        child states, produce a compressed parent state that subsumes
        the children's information.
  \item $\info: \States \to \Real_{\geq 0}$ is the \emph{information
        metric}: a measure of the content in a state.
\end{itemize}
\end{definition}

We also explore an optional Hilbert locality optimization (Appendix~\ref{app:hilbert}), though ablation experiments show no measurable benefit (Section~\ref{sec:ablation}).

\begin{definition}[Information Preservation Ratio]
\label{def:preservation}
The \emph{preservation ratio} of a summarization function $\summarize$
under metric $\info$ is:
\begin{equation}
  \preserv(\summarize, \info) \;=\;
  \inf_{s_p, s_{c_1}, \dots, s_{c_k}} \;
  \frac{\info\bigl(\summarize(s_p, [s_{c_1}, \dots, s_{c_k}])\bigr)}
       {\info(s_p) + \sum_{i=1}^{k} \info(s_{c_i})}
  \label{eq:preservation}
\end{equation}
We require $\preserv \geq 1 - \epsilon$ for a configurable $\epsilon > 0$.
\end{definition}


% ============================================================
\subsection{Architecture}
\label{sec:architecture}

The BIC exposes four operations: \textsc{Spawn}, \textsc{Execute},
\textsc{Query}, and \textsc{Terminate}.  Figure~\ref{fig:architecture}
shows the data flow.

\begin{figure}[t]
\centering
\begin{tikzpicture}[
  block/.style={draw, rounded corners, minimum width=2.4cm,
                minimum height=0.8cm, font=\small},
  arr/.style={-{Stealth[length=2mm]}, thick},
  every node/.style={font=\small}
]
  % Blocks — wider horizontal spacing to avoid overlaps
  \node[block, fill=blue!10] (api) {API Layer};
  \node[block, fill=green!10, below left=1.2cm and 1.4cm of api] (reg) {Agent Registry};
  \node[block, fill=orange!10, below right=1.2cm and 1.4cm of api] (cache) {BIC (Fixed Array)};
  \node[block, fill=purple!10, below=1.2cm of reg] (cantor) {Cantor Addr};
  \node[block, fill=red!10, below=1.2cm of cache] (evict) {Eviction Mgr};
  \node[block, fill=cyan!10, left=1.2cm of evict] (hilbert) {Hilbert Idx};

  % Arrows — re-routed to avoid crossing over blocks
  \draw[arr] (api) -- node[left, pos=0.4] {register} (reg);
  \draw[arr] (api) -- node[right, pos=0.4] {put/get} (cache);
  \draw[arr] (reg) -- node[above, pos=0.4] {slot hint} (cache);
  \draw[arr] (reg) -- node[left] {$\alpha(v)$} (cantor);
  \draw[arr] (cache) -- node[right] {full?} (evict);
  \draw[arr] (evict) -- node[above] {cluster} (hilbert);
  \draw[arr] (evict.north) -- node[left] {$\summarize$} (cache.south);
\end{tikzpicture}
\caption{BIC architecture.  The API layer coordinates the Agent
  Registry (tree tracking + Cantor addressing), the fixed-size cache
  array, and the Eviction Manager (hierarchical summarization, with
  optional Hilbert-locality clustering).}
\label{fig:architecture}
\end{figure}

Algorithm~\ref{alg:spawn} shows the \textsc{Spawn} operation;
Algorithm~\ref{alg:query} shows \textsc{Query} with
ancestor-chain reconstruction.

\begin{algorithm}[t]
\caption{\textsc{Spawn}$(v_{\mathrm{parent}}, \tau)$ — spawn a new
  agent with task $\tau$ as a child of $v_{\mathrm{parent}}$.}
\label{alg:spawn}
\begin{algorithmic}[1]
  \STATE $v \gets \mathrm{Registry.Register}(v_{\mathrm{parent}}, \tau)$
  \COMMENT{Heuristic, no LLM: assign tree position, compute Cantor addr.}
  \IF{$\mathrm{Cache.IsFull}()$}
    \STATE $\mathrm{EvictionManager.EnsureCapacity}(1)$
    \COMMENT{Triggers Algorithm~\ref{alg:evict}; summarization is LLM-eligible}
  \ENDIF
  \STATE $s_0 \gets \{\texttt{task}: \tau\}$
  \COMMENT{Heuristic, no LLM}
  \STATE $\eta \gets \hilbert(\mathrm{ExtractVector}(s_0))$
  \COMMENT{Closed-form bit-interleaving (Butz~1971), no learning}
  \STATE $\mathrm{Cache.Put}(v, s_0, \eta, h(v))$
  \COMMENT{$h(v) = \cantor(v) \bmod \slots$; closed-form (Cantor 1878)}
  \RETURN $v$
\end{algorithmic}
\end{algorithm}

\begin{algorithm}[t]
\caption{\textsc{Query}$(v)$ — retrieve agent $v$'s state.}
\label{alg:query}
\begin{algorithmic}[1]
  \STATE $e \gets \mathrm{Cache.Get}(v)$
  \COMMENT{Heuristic Cantor-index lookup, no LLM}
  \IF{$e \neq \mathrm{nil}$}
    \RETURN $e.\mathrm{state}$
    \COMMENT{Cache hit: $\Oh(1)$, no LLM}
  \ENDIF
  \STATE $r \gets \mathrm{Registry.Get}(v)$
  \IF{$r = \mathrm{nil}$}
    \RETURN nil
  \ENDIF
  \STATE \COMMENT{Reconstruct from ancestor chain: $\Oh(d)$ — heuristic walk}
  \STATE $\mathrm{summaries} \gets [\,]$
  \STATE $u \gets r.\mathrm{parent}$
  \WHILE{$u \neq \mathrm{nil}$}
    \STATE $e_u \gets \mathrm{Cache.Get}(u)$
    \IF{$e_u \neq \mathrm{nil}$}
      \STATE Append $e_u.\mathrm{state.\_summaries}$ to summaries
      \STATE \textbf{break}
      \COMMENT{Found cached ancestor; LLM optional for re-expansion}
    \ENDIF
    \STATE $u \gets \mathrm{Registry.Get}(u).\mathrm{parent}$
  \ENDWHILE
  \RETURN $\mathrm{MergeReconstruction}(r, \mathrm{summaries})$
  \COMMENT{Heuristic merge by default; LLM-eligible re-expansion if equipped}
\end{algorithmic}
\end{algorithm}

\paragraph{Eviction policy.}
When the cache is full, the eviction manager selects victims using a
three-level priority: (1)~idle agents (not accessed for $>t$ seconds);
(2)~deepest agents (largest tree depth); (3)~least recently accessed.
Selected agents are evicted and their states are summarized into their
respective parents via $\summarize$.  An optional Hilbert-locality
clustering step is described in Appendix~\ref{app:hilbert}.

\paragraph{What is heuristic vs.\ LLM-based?}
To make the implementation choice explicit, Table~\ref{tab:mechanisms}
maps each step of the BIC pipeline to the concrete mechanism that
implements it.  The addressing layer (Cantor and Hilbert) is entirely
closed-form arithmetic with no learning component; the eviction
selector is a deterministic three-level priority queue; the
summarization step is the only LLM-eligible component, and our
default is a deterministic concatenating summarizer (preserves
recency, no LLM call) with a pluggable interface for an LLM
summarizer (Gemini-2.0-flash in our experiments) using the prompt
template in Appendix~\ref{app:summary-prompt}.

\begin{table}[t]
\centering
\small
\caption{Each BIC step and the mechanism implementing it.  Only
  hierarchical summarization is LLM-eligible; all addressing and
  selection logic is closed-form.}
\label{tab:mechanisms}
\vspace{0.3em}
\begin{tabular}{lll}
\toprule
\textbf{Step} & \textbf{Mechanism} & \textbf{LLM?} \\
\midrule
Cantor pairing $\cantor$ & Closed-form arithmetic (Cantor~1878) & No \\
Hilbert index $\hilbert$ & Bit-interleaving (Butz~1971) & No \\
Slot mapping $h(v) = \cantor(v) \bmod \slots$ & Modular arithmetic & No \\
Spawn (\textsc{Register}) & Heuristic registry write & No \\
Eviction priority & Three-level heuristic (idle / depth / LRU) & No \\
Hilbert clustering (optional) & Sort + sliding-window selection & No \\
Hierarchical summarization $\summarize$ & Default: concat-and-merge & No \\
\quad\emph{LLM-equipped variant} & Gemini-2.0-flash via plug-in & Yes \\
Query (cache hit) & $\Oh(1)$ hash lookup & No \\
Query (reconstruct) & Ancestor-chain walk + merge & No (default) \\
\quad\emph{LLM re-expansion} & Gemini-2.0-flash on cached summaries & Yes \\
\bottomrule
\end{tabular}
\end{table}


% ============================================================
\subsection{Theorems}
\label{sec:theorems}

We state four theorems; full proofs are in Appendix~\ref{sec:proofs}.

\begin{theorem}[Bounded Memory]
\label{thm:bounded}
For any execution trace $\tau$ of length $|\tau| \to \infty$, the BIC
uses at most $\slots \cdot C_{\max} + \Oh(N)$ total memory, where
$C_{\max} = \max_{s \in \States} |s|$ is the maximum state size and
$N$ is the total number of agents ever spawned.  The $\Oh(N)$ term
accounts for the agent registry (lightweight metadata: parent pointer,
depth, status --- $\Oh(1)$ per agent).
\end{theorem}

\begin{remark}
The registry's $\Oh(N)$ growth is an honest limitation.  Each registry
record is $\sim$100 bytes (agent ID, parent ID, depth, status),
compared to 1--4\,KB per agent state.  At $N = 10{,}000$ agents, the
registry uses $\sim$1\,MB vs.\ $\sim$40\,MB for a cache of
$\slots=1{,}000$ states.  In practice, the registry can also be bounded
with a secondary eviction policy; we do not do so here to keep the
architecture simple.
\end{remark}

\begin{theorem}[Fragmentation]
\label{thm:frag}
After a call to \textsc{Compact}, the BIC has zero external
fragmentation: all $|\BIC|$ occupied slots are contiguous in
positions $[0, |\BIC|)$.
\end{theorem}

\begin{theorem}[Retrievability]
\label{thm:retrieve}
Every agent state is retrievable:
\begin{enumerate}
  \item[(a)] \textbf{Cached agents:} Retrieved in $\Oh(1)$ via hash-map
    lookup.
  \item[(b)] \textbf{Evicted agents:} Reconstructed in $\Oh(d)$ by
    walking the ancestor chain (depth $d$), with information
    loss bounded by $1 - \preserv^d$ where $\preserv \geq 1-\epsilon$
    is the preservation ratio (Definition~\ref{def:preservation}).
    For branching factor $k$, $d \leq \log_k N$, so
    reconstruction is $\Oh(\log_k N)$.
\end{enumerate}
\end{theorem}

\begin{theorem}[Indefinite Liveness]
\label{thm:liveness}
The BIC system can execute indefinitely: for any spawn event,
\textsc{EnsureCapacity} either finds a free slot or successfully
evicts $\geq 1$ agent.  Every operation completes in $\Oh(\slots)$
worst-case time, independent of $N$.
\end{theorem}
